\chapter{CƠ SỞ LÝ THUYẾT VÀ CÁC CÔNG TRÌNH LIÊN QUAN}
\label{chap:chapter2}

\section{Lý thuyết kiểm soát tồn kho ngẫu nhiên}
\label{sec:lyThuyetTonKho}

\subsection{Mô hình tồn kho và các thành phần chi phí}
\label{subsec:moHinhTonKho}

Bài toán kiểm soát tồn kho xoay quanh việc xác định thời điểm và số lượng đặt
hàng bổ sung nhằm cân bằng giữa chi phí phát sinh khi dự trữ quá nhiều và rủi
ro khi dự trữ quá ít \parencite{zipkin2000}. Xét một hệ thống đơn kho -- đơn
mặt hàng hoạt động theo chu kỳ rời rạc. Gọi $I_t$ là mức tồn kho tại đầu kỳ
$t$, $a_t$ là số lượng đặt hàng, $D_t$ là nhu cầu phát sinh trong kỳ $t$ (biến
ngẫu nhiên) và $L_t$ là thời gian giao hàng. Chi phí vận hành mỗi kỳ bao gồm:

\begin{itemize}
    \item Chi phí lưu kho $C_{\text{hold}} = c_{lk} \cdot \max(I_t, 0)$:
    chi phí tỷ lệ với lượng hàng tồn dương trong kho, bao gồm chi phí vốn,
    bảo quản và kho bãi.

    \item Chi phí thiếu hàng $C_{\text{stock}} = c_{th} \cdot \max(D_t - I_t, 0)$:
    chi phí phát sinh khi nhu cầu vượt quá lượng hàng có sẵn. Trong mô hình
    mất đơn hàng (lost sales), phần nhu cầu không được đáp ứng bị mất hoàn toàn
    thay vì chuyển sang kỳ sau \parencite{zipkin2008}.

    \item Chi phí đặt hàng cố định $C_{\text{order}} = c_{dh} \cdot \mathbb{1}(a_t > 0)$:
    chi phí phát sinh mỗi khi hệ thống phát lệnh đặt hàng, bất kể số lượng đặt.

    \item Phạt tràn kho $C_{\text{over}} = p_{tk} \cdot \max(I_t + q_t - W, 0)$:
    chi phí phát sinh khi tổng lượng hàng tồn và hàng nhập vượt quá sức chứa
    $W$ của kho, trong đó $q_t$ là lượng hàng nhập kho trong kỳ.

    \item Phạt vi phạm mức phục vụ dựa trên tỷ lệ đáp ứng nhu cầu
    (fill rate) trên cửa sổ trượt, phản ánh yêu cầu kinh doanh về chất lượng
    dịch vụ khách hàng.
\end{itemize}

Tổng chi phí vận hành của một cặp kho--mặt hàng trong kỳ $t$ được tính:
\begin{equation}
    C_t = C_{\text{hold}} + C_{\text{stock}} + C_{\text{order}} + C_{\text{over}}
    + \phi_{dv} \cdot \max\bigl(\tau - \text{FR}_t^{(w)},\; 0\bigr)
    \label{eq:tongchiphi}
\end{equation}
trong đó $\text{FR}_t^{(w)}$ là tỷ lệ đáp ứng trên cửa sổ trượt $w$ ngày,
$\tau$ là ngưỡng mức phục vụ mục tiêu và $\phi_{dv}$ là hệ số phạt.

\subsection{Các chính sách đặt hàng lại cổ điển}
\label{subsec:chinhSachCoDien}

Ba chính sách cổ điển được sử dụng làm phương pháp đối chứng trong khóa luận.

Lượng đặt hàng kinh tế (EOQ). Harris \cite{harris1913} đề xuất công thức
tối ưu hóa tần suất đặt hàng dưới giả định nhu cầu xác định và đều:
\begin{equation}
    Q^* = \sqrt{\frac{2 D \cdot c_{dh}}{c_{lk}}}
    \label{eq:eoq}
\end{equation}
trong đó $D$ là nhu cầu trung bình hàng năm. Mô hình giả định không có thiếu
hàng, thời gian giao hàng cố định và không có ràng buộc sức chứa.

Chính sách $(s, S)$. Scarf \cite{scarf1960} chứng minh rằng chính sách
dạng $(s, S)$ --- đặt hàng để nâng mức tồn kho lên $S$ khi mức hiện tại giảm
xuống dưới $s$ --- là tối ưu cho hệ đơn kho với nhu cầu ngẫu nhiên, chi phí đặt
hàng cố định, và cơ chế giữ đơn hàng chờ (backorder). Ngưỡng $s$ thường được
tính dựa trên nhu cầu trung bình và độ lệch chuẩn trong thời gian giao hàng:
\begin{equation}
    s = \mu_L + z_\alpha \cdot \sigma_L
    \label{eq:sSPolicy}
\end{equation}
trong đó $\mu_L$ và $\sigma_L$ là kỳ vọng và độ lệch chuẩn của nhu cầu tích
lũy trong thời gian giao hàng, $z_\alpha$ là phân vị chuẩn tương ứng với mức
phục vụ mong muốn.

Mô hình Newsvendor. Arrow và cộng sự \cite{arrow1951} đề xuất mô hình cho bài toán
đặt hàng đơn kỳ, trong đó lượng đặt hàng tối ưu là:
\begin{equation}
    Q^* = F^{-1}\!\left(\frac{c_{th}}{c_{th} + c_{lk}}\right)
    \label{eq:newsvendor}
\end{equation}
với $F^{-1}$ là hàm phân vị của phân phối nhu cầu. Mô hình chỉ tối ưu cho một
chu kỳ duy nhất, không xét lượng hàng đang trên đường vận chuyển.

Cả ba chính sách trên đều được xây dựng cho bài toán đơn kho -- đơn mặt hàng
và không tối ưu hóa đồng thời ràng buộc sức chứa dùng chung, phạt tràn kho và
mức phục vụ --- những đặc trưng quan trọng trong bài toán đa kho -- đa mặt
hàng mà khóa luận xét \parencite{zipkin2000, silver2016}.

\subsection{Các thước đo mức phục vụ trong quản lý tồn kho}
\label{subsec:thuocDoMucPhucVu}

Trong quản lý tồn kho ngẫu nhiên, mức phục vụ (service level) được đánh giá thông qua ba thước đo phổ biến \parencite{silver2016, zipkin2000}:
\begin{itemize}
    \item \textit{Mức phục vụ theo chu kỳ} ($\alpha$ -- Cycle Service Level / Type-1): xác suất không xảy ra hiện tượng thiếu hàng trong một chu kỳ bổ sung tồn kho, tức $P(D_L \le s)$ trong đó $D_L$ là nhu cầu tích lũy trong thời gian giao hàng $L$ và $s$ là điểm đặt hàng lại. Tỷ số tới hạn trong mô hình Newsvendor (\ref{eq:newsvendor}) phản ánh chính xác phân vị $z_\alpha$ của thước đo này.
    \item \textit{Tỷ lệ đáp ứng nhu cầu} ($\beta$ -- Fill Rate / Type-2): tỷ lệ phần trăm tổng nhu cầu của khách hàng được đáp ứng ngay lập tức từ lượng hàng sẵn có trong kho.
    \item \textit{Tỷ lệ sẵn sàng} ($\gamma$ -- Ready Rate): tỷ lệ thời gian mà mức tồn kho khả dụng duy trì giá trị dương.
\end{itemize}

Cần lưu ý rạch ròi rằng tỷ số chi phí $c_{th}/(c_{th} + c_{lk})$ xác định mức phục vụ chu kỳ $\alpha$, chứ không phải tỷ lệ fill rate $\beta$. Trong điều kiện nhu cầu ổn định và biến động vừa phải, $\beta$ thường cao hơn $\alpha$. Tuy nhiên, dưới tác động của nhu cầu gián đoạn (intermittent demand) và thời gian giao hàng ngẫu nhiên như trong dữ liệu M5 Walmart, mối quan hệ này không còn mang tính tuyệt đối và $\beta$ có thể sụt giảm thấp hơn $\alpha$ \parencite{silver2016}.

Một quan sát thực nghiệm liên quan: với bộ tham số chi phí $c_{th} = 10{,}0$ và
$c_{lk} = 1{,}0$ sử dụng trong thực nghiệm, tỷ số tới hạn lý thuyết của Newsvendor
là $c_{th}/(c_{th}+c_{lk}) = 10/11 \approx 0{,}909$. Tuy nhiên, tìm kiếm lưới
trên môi trường mô phỏng cho tham số tối ưu $CR^* = 0{,}85$ --- thấp hơn nhưng
không lệch quá xa lý thuyết. Chênh lệch còn lại không phải lỗi cài đặt mà phản
ánh sự khác biệt khái niệm
căn bản: tỷ số tới hạn cổ điển tối ưu cho mức phục vụ chu kỳ $\alpha$ trong bài
toán \textit{một kỳ} không có hàng đang vận chuyển, còn môi trường thực nghiệm
là bài toán \textit{đa kỳ} với pipeline $3$ ngày, ràng buộc sức chứa kho và hàm
mục tiêu fill rate trên cửa sổ trượt. Khoảng cách này được ghi nhận và phân tích
trong mục thảo luận kết quả (Mục~\ref{subsec:isoService}).

\section{Nền tảng học tăng cường}
\label{sec:nenTangRL}

\subsection{Quá trình quyết định Markov (MDP)}
\label{subsec:mdp}

Bài toán ra quyết định tuần tự dưới bất định được mô hình hóa bằng quá trình
quyết định Markov \parencite{puterman1994}, xác định bởi bộ năm
$(S, A, P, R, \gamma)$ trong đó:
\begin{itemize}
    \item $S$: không gian trạng thái (state space).
    \item $A$: không gian hành động (action space).
    \item $P(s' \mid s, a)$: hàm xác suất chuyển trạng thái.
    \item $R(s, a)$: hàm phần thưởng (reward function).
    \item $\gamma \in [0, 1)$: hệ số chiết khấu (discount factor).
\end{itemize}

Mục tiêu của tác tử là tìm chính sách $\pi^*$ cực đại hóa lợi tức kỳ vọng:
\begin{equation}
    \pi^* = \arg\max_\pi \; \mathbb{E}_\pi \left[ \sum_{t=0}^{\infty}
    \gamma^t R(s_t, a_t) \right]
    \label{eq:objectiveRL}
\end{equation}

\subsection{Hàm giá trị và phương trình Bellman}
\label{subsec:hamGiaTri}

Hàm giá trị trạng thái $V^\pi(s) = \mathbb{E}_\pi[\sum_{t=0}^{\infty} \gamma^t
R_t \mid s_0 = s]$ và hàm giá trị hành động $Q^\pi(s, a)$ thỏa mãn phương trình
Bellman \parencite{sutton2018}. Hàm lợi thế (advantage function) được định nghĩa:
\begin{equation}
    A^\pi(s, a) = Q^\pi(s, a) - V^\pi(s)
    \label{eq:advantage}
\end{equation}
đo mức độ tốt hơn (hoặc kém hơn) của hành động $a$ so với hành vi trung bình
của chính sách $\pi$ tại trạng thái $s$.

\subsection{Ước lượng lợi thế tổng quát (GAE)}
\label{subsec:gae}

Schulman và cộng sự \cite{schulman2016gae} đề xuất phương pháp ước lượng lợi thế tổng quát
(Generalized Advantage Estimation -- GAE) nhằm kiểm soát đánh đổi giữa thiên
lệch (bias) và phương sai (variance) trong ước lượng gradient chính sách:
\begin{equation}
    \hat{A}_t^{\text{GAE}(\gamma, \lambda)} = \sum_{l=0}^{\infty}
    (\gamma \lambda)^l \delta_{t+l}
    \label{eq:gae}
\end{equation}
trong đó $\delta_t = r_t + \gamma V(s_{t+1}) - V(s_t)$ là sai số TD (temporal
difference) và $\lambda \in [0, 1]$ là tham số kiểm soát đánh đổi. Khi
$\lambda = 0$, GAE suy biến thành ước lượng một bước (thiên lệch cao, phương sai
thấp); khi $\lambda = 1$, GAE tương đương Monte-Carlo (không thiên lệch, phương
sai cao). Khóa luận sử dụng $\lambda = 0{,}95$ theo khuyến nghị thực nghiệm
phổ biến.

\subsection{Thuật toán PPO (Proximal Policy Optimization)}
\label{subsec:ppo}

Schulman và cộng sự \cite{schulman2017ppo} đề xuất PPO nhằm khắc phục nhược điểm của phương
pháp tối ưu hóa chính sách: các bước cập nhật quá lớn có thể phá hủy chính
sách đã học. PPO sử dụng hàm mục tiêu cắt (clipped surrogate objective):
\begin{equation}
    L^{\text{CLIP}}(\theta) = \mathbb{E}_t \left[ \min\!\left(
    r_t(\theta) \hat{A}_t, \;
    \text{clip}\bigl(r_t(\theta), 1 - \epsilon, 1 + \epsilon\bigr) \hat{A}_t
    \right) \right]
    \label{eq:ppoclip}
\end{equation}
trong đó $r_t(\theta) = \frac{\pi_\theta(a_t \mid s_t)}{\pi_{\theta_{\text{old}}}(a_t \mid s_t)}$
là tỷ lệ xác suất giữa chính sách mới và cũ, và $\epsilon$ (thường bằng
$0{,}2$) giới hạn biên độ thay đổi tối đa. PPO kết hợp cập nhật on-policy, hiệu
quả tính toán cao và tính ổn định, phù hợp cho bài toán có không gian hành động
rời rạc hoặc liên tục.

Ý tưởng cốt lõi của TRPO \parencite{schulman2015trpo} --- giới hạn bước cập
nhật trong vùng tin cậy --- được PPO thực hiện đơn giản hơn qua phép cắt, giúp
cài đặt dễ dàng và huấn luyện nhanh hơn trên thực tế.

\section{Học tăng cường đa tác tử}
\label{sec:marl}

\subsection{Mô hình Dec-POMDP}
\label{subsec:decPomdp}

Khi hệ thống có nhiều tác tử ra quyết định đồng thời, bài toán được mô hình hóa
bằng quá trình quyết định Markov phi tập trung có quan sát bộ phận (Decentralized
Partially Observable Markov Decision Process -- Dec-POMDP)
\parencite{oliehoek2016}. Trong mô hình này, mỗi tác tử $i$ chỉ quan sát được
một phần trạng thái toàn cục thông qua hàm quan sát $o_i = O(s, i)$ và ra quyết
định dựa trên quan sát cục bộ. Bài toán tìm chính sách chung tối ưu cho tất cả
tác tử đã được chứng minh thuộc lớp NEXP-đầy đủ \parencite{bernstein2002},
nên các lời giải chính xác không khả thi ở quy mô lớn.

Bài toán tồn kho đa kho đa mặt hàng được mô hình hóa tự nhiên dưới dạng Dec-POMDP
vì bốn đặc điểm cấu trúc phù hợp với lớp bài toán này:
\begin{enumerate}
    \item \textbf{Quyết định tuần tự, hệ quả trễ}. Đơn đặt hàng hôm nay ảnh hưởng
    tồn kho $1$--$3$ ngày sau. Không phải bài toán một bước.
    \item \textbf{Trạng thái thay đổi theo hành động}. Đặt hàng làm đổi tồn
    kho, thay đổi luôn bài toán của ngày mai.
    \item \textbf{Không có nhãn đúng}. Không ai biết lượng đặt tối ưu cho ngày
    $t$ là bao nhiêu; chỉ biết chi phí sau khi thực hiện. Đây là lý do không
    dùng học có giám sát.
    \item \textbf{Mục tiêu dài hạn}. Tối thiểu chi phí cả năm, không phải
    chỉ hôm nay. Tích trữ hôm nay tốn tiền nhưng tránh mất doanh thu ngày mai.
\end{enumerate}

Quy hoạch động không khả dụng vì không gian trạng thái liên tục $43$ chiều
$\times\,300$ tác tử và phân phối nhu cầu M5 không có dạng đóng --- đây
là ``lời nguyền số chiều'' mà các phương pháp bảng tra (tabular) không vượt qua
được \parencite{bellman1957}.

\subsection{Học độc lập (Independent Learning)}
\label{subsec:hocDocLap}

Tan \cite{tan1993} đề xuất phương pháp học độc lập: mỗi tác tử coi các tác tử
khác là một phần của môi trường và áp dụng thuật toán RL đơn tác tử riêng rẽ.
Hạn chế lý thuyết chính là môi trường trở nên không dừng (non-stationary) do các
tác tử khác cũng đang thay đổi chính sách. Tuy nhiên, các nghiên cứu thực
nghiệm cho thấy phương pháp này đạt hiệu quả bất ngờ trên nhiều bài toán hợp
tác \parencite{dewitt2020ippo}.

\subsection{IPPO — Independent Multi-Agent PPO}
\label{subsec:ippo}

IPPO là biến thể đa tác tử của PPO trong đó mỗi tác tử $i$ chạy PPO độc lập
dựa trên quan sát cục bộ $o_i^t$ và phần thưởng riêng $r_i^t$
\parencite{dewitt2020ippo}. Yu và cộng sự \cite{yu2022mappo} đánh giá hệ thống trên nhiều
môi trường hợp tác và kết luận rằng IPPO đạt hiệu năng cạnh tranh hoặc vượt
trội so với các phương pháp phức tạp hơn như QMIX hay MAPPO trong nhiều kịch bản.

Đặc điểm quan trọng của IPPO khiến nó phù hợp cho bài toán tồn kho:
\begin{itemize}
    \item Không yêu cầu trạng thái toàn cục: mỗi tác tử (cặp kho--mặt hàng)
    chỉ cần quan sát cục bộ.
    \item Khả năng mở rộng: không gian hành động của mỗi tác tử có kích thước
    $K$ cố định, không phụ thuộc số tác tử.
    \item Tính đơn giản: cài đặt dựa trên PPO chuẩn, không cần mixing network
    hay cơ chế truyền thông.
\end{itemize}

Cần phân biệt rõ ba tầng nghĩa của ``independent'' trong IPPO, vì chúng
hàm ý khác nhau:
\begin{itemize}
    \item \textit{Độc lập về ra quyết định}: mỗi tác tử chọn hành động chỉ dựa
    trên quan sát cục bộ của mình, không nhận thông tin từ tác tử khác. Đúng.
    \item \textit{Độc lập về tính lợi thế}: GAE tính riêng theo trục thời gian
    cho từng cặp $i$, sử dụng phần thưởng riêng $r_i^t$ và giá trị riêng $V(o_i^t)$.
    Đúng.
    \item \textit{Độc lập về tham số}: mỗi tác tử có bộ trọng số riêng. Không
    --- đây là điểm chia sẻ tham số (parameter sharing) được khảo sát ở
    Mục~\ref{subsec:paramSharing}.
\end{itemize}

Hai cặp ở hai trạng thái khác nhau sẽ hành động khác nhau dù cùng $\theta$
--- vì đầu vào khác nhau. Mỗi tác tử có chính sách riêng theo nghĩa hàm
$\pi(a \mid o_i)$ nhận quan sát riêng và cho ra phân phối riêng, nhưng
tham số $\theta$ dùng chung.

\subsection{Chia sẻ tham số (Parameter Sharing)}
\label{subsec:paramSharing}

Gupta và cộng sự \cite{gupta2017} chỉ ra rằng khi các tác tử có cùng không gian quan sát và
hành động, một mạng nơ-ron duy nhất có thể được chia sẻ giữa tất cả tác tử.
Mỗi tác tử vẫn ra quyết định dựa trên quan sát riêng nhưng dùng chung bộ trọng
số. Lợi ích:
\begin{itemize}
    \item Số tham số cần học không tăng theo số tác tử (ở đây: $300$).
    \item Mỗi bước cập nhật sử dụng dữ liệu từ tất cả tác tử, tăng hiệu quả mẫu.
    \item Chính sách học được có khả năng tổng quát hóa cho các mặt hàng có quy
    mô nhu cầu khác nhau.
\end{itemize}

Chẩn đáng lưu ý: chia sẻ tham số không tương đương với việc gộp $300$ tác tử
thành một tác tử duy nhất. Nếu gộp, không gian hành động sẽ là $6^{300}$,
không thể học được trên thực tế. Với kiến trúc chia sẻ tham số hiện tại,
mạng nhận batch $(300, 43)$ và forward một lần, nhưng mỗi tác tử chỉ xử lý
$K = 6$ hành động --- không gian cố định và có thể giải quyết được.

\subsection{Phần thưởng phân rã cục bộ và bài toán gán tín dụng}
\label{subsec:rewardDecomposition}

Trong hệ đa tác tử, việc sử dụng phần thưởng toàn cục (tổng chi phí của toàn bộ
hệ thống) cho mọi tác tử gây ra vấn đề gán tín dụng (credit assignment): tác tử
khó phân biệt được đóng góp cá nhân trong tín hiệu tổng hợp
\parencite{wolpert2001}. Khóa luận áp dụng phân rã cục bộ: mỗi tác tử nhận
phần thưởng riêng dựa trên chi phí và tỷ lệ đáp ứng của chính cặp kho--mặt
hàng mà nó quản lý.

Khung lý thuyết định hình phần thưởng bảo toàn thế năng (Potential-Based Reward Shaping -- PBRS) do Ng và cộng sự \cite{ng1999} chứng minh khẳng định rằng: một hàm biến đổi phần thưởng $F(s,a,s')$ bảo toàn tập chính sách tối ưu của MDP gốc khi và chỉ khi nó có dạng chênh lệch thế năng $F(s,a,s') = \gamma\Phi(s') - \Phi(s)$. Khung lý thuyết này sau đó được mở rộng sang hệ đa tác tử bởi Devlin và Kudenko \cite{devlin2011potential} cũng như ứng dụng trong tồn kho bởi De Moor và cộng sự \cite{demoor2022reward}.

\section{Các công trình nghiên cứu liên quan}
\label{sec:congTrinh}

\subsection{Học tăng cường trong quản lý tồn kho}
\label{subsec:rlTonKho}

Oroojlooyjadid và cộng sự \cite{oroojlooyjadid2022} áp dụng Deep Q-Network (DQN) cho trò chơi phân phối bia (beer game) --- một bài toán chuỗi cung ứng 4 cấp cổ điển --- và cho thấy DQN có thể học được chính sách giảm hiệu ứng bullwhip. Tuy nhiên, nghiên cứu giới hạn ở hệ đơn mặt hàng, bốn tác tử.

Gijsbrechts và cộng sự \cite{gijsbrechts2022} đánh giá hệ thống khả năng của DRL trên ba lớp bài toán tồn kho: mất đơn hàng, hai nguồn cung và đa cấp. Kết quả cho thấy DRL đạt hiệu năng ngang hoặc tốt hơn chính sách cổ điển trong hầu hết trường hợp, nhưng tất cả thực nghiệm đều ở quy mô nhỏ (dưới 10 SKU).

Hubbs và cộng sự \cite{hubbs2020orgym} giới thiệu bộ thư viện OR-Gym đóng vai trò là môi trường benchmark chuẩn cho các bài toán Vận trù học bằng RL, bao gồm cả mô hình Newsvendor và quản lý tồn kho đa cấp lost-sales.

Nghiên cứu của Ding và cộng sự \cite{ding2022cdppo} (CD-PPO) phát biểu chính xác bài toán tồn kho đơn cửa hàng nhiều SKU dưới dạng tài nguyên kho dùng chung làm ghép nối các tác tử vốn độc lập. Họ so sánh trực tiếp IPPO và MAPPO có chia sẻ tham số (parameter sharing), chứng minh hiệu quả vượt trội khi quản lý tài nguyên kho chung.

Yang và cộng sự \cite{yang2023mabim} phát triển MABIM --- bộ benchmark MARL đa dạng cho tồn kho. MABIM coi mỗi SKU tại mỗi điểm lưu trữ là một tác tử độc lập, sử dụng dữ liệu nhu cầu thực tế của hơn 2.000 SKU với 51 tác vụ thử nghiệm (bao gồm nhóm tác vụ mở rộng quy mô scaling up), thử nghiệm các thuật toán MARL bao gồm IPPO dưới ràng buộc sức chứa dùng chung.

Kotecha và del Rio Chanona \cite{kotecha2024decentralized} tiến hành phân tích chuyên sâu các thuật toán MARL cho kiểm soát tồn kho phân tán, so sánh giữa IPPO thuần, IPPO chia sẻ tham số và MAPPO. Họ phát hiện rằng mặc dù parameter sharing giảm số lượng tham số cần học, hiệu năng của IPPO chia sẻ tham số có dấu hiệu giảm sút khi quy mô tác tử tăng lên rất lớn.

Gần đây nhất, Liu và cộng sự \cite{liu2025happo} áp dụng thuật toán Heterogeneous-Agent PPO (HAPPO) cho tồn kho đa cấp. Đáng chú ý, các tác giả phát hiện việc kết hợp chi phí cục bộ và chi phí toàn hệ thống vào tín hiệu thưởng thu được chính sách có độ ổn định và hiệu năng cao hơn so với chỉ dùng chi phí toàn hệ thống thuần túy.

\subsection{Hai công trình gần nhất cùng bối cảnh}
\label{subsec:congTrinhGanNhat}

Hai công trình công bố gần đây có bối cảnh gần với khóa luận nhất và vì vậy cần
được đối chiếu chi tiết.

Zhu và Wu \cite{zhu2026spatiotemporal} đề xuất khung DC-STGN-MARL kết hợp một
mạng đồ thị không--thời gian cho dự báo nhu cầu với học tăng cường đa tác tử
theo hướng phân rã hàm giá trị, đánh giá trên chính bộ dữ liệu M5. Công trình
báo cáo mức giảm $13{,}6\%$ tổng chi phí so với học tăng cường đơn tác tử và
mức phục vụ $95{,}8\%$. Điểm cần lưu ý đối với khóa luận nằm ở khảo sát loại
trừ của họ: khi thay cơ chế phối hợp bằng học độc lập, chi phí tăng từ $125{,}4$
lên $152{,}7$ và mức phục vụ giảm từ $95{,}8\%$ xuống $92{,}1\%$. Kết quả này
đặt ra câu hỏi trực tiếp về lựa chọn học độc lập của khóa luận.

Tuy nhiên, bối cảnh của hai nghiên cứu khác nhau ở một điểm mấu chốt. Môi
trường của Zhu và Wu giả định năng lực cung ứng không giới hạn và không đưa sức
chứa kho vào hàm mục tiêu, nên các tác tử chỉ ghép nối với nhau qua hàm phần
thưởng chung chứ không qua một nguồn lực vật lý hữu hạn. Trong khóa luận, sức
chứa kho dùng chung là một ràng buộc cứng và cũng chính là kênh phối hợp gián
tiếp (Mục~\ref{subsec:phoiHop}); mức độ cần thiết của phối hợp tường minh trong
bối cảnh này vì vậy là một câu hỏi thực nghiệm mở, chứ không thể suy trực tiếp
từ kết quả của họ. Ngoài ra, một số chi tiết trong công trình cần được tiếp
nhận thận trọng: các so sánh chi phí được thực hiện giữa những cấu hình có mức
phục vụ khác nhau, và số lượng tác tử không được xác định rõ ràng --- bài viết
mô tả $N$ tác tử ``đại diện cho $N$ kho hoặc nhóm sản phẩm'' mà không xác định cụ thể.

Ở hướng thứ hai, một khung làm việc dựa trên kiến trúc hướng sự kiện kết hợp
tác tử xếp vị trí lưu trữ và tác tử bổ sung hàng đã được đề xuất
\cite{jeong2026alca}, với đóng góp chính là cơ chế \textit{gán công liên kết
hành động}: phần thưởng chỉ được tính khi lô hàng thực nhập kho sau thời gian
giao hàng, rồi được quy ngược về đúng cặp trạng thái--hành động đã phát lệnh
đặt. Đây là một cách xử lý trực diện độ trễ giữa quyết định và hệ quả của nó.
Khảo sát loại trừ của công trình cho thấy khi bỏ cơ chế này, tỷ lệ thiếu hàng
tăng vọt. Cần lưu ý rằng biến thể đối chứng của họ sử dụng trạng thái
dự đoán tại thời điểm ra quyết định, chứ không phải ước lượng lợi thế
tổng quát như khóa luận sử dụng; do đó con số của họ không thể chuyển trực tiếp
sang bối cảnh này. Hạn chế của công trình là chỉ xét một kho duy nhất, trên dữ
liệu tổng hợp, và hàm phần thưởng không quy về đơn vị tiền tệ.

Bảng~\ref{tab:soSanhCongTrinh} tóm tắt đối chiếu.

\begin{table}[htbp]
    \centering
    \caption{Đối chiếu khóa luận với hai công trình gần nhất cùng bối cảnh}
    \label{tab:soSanhCongTrinh}
    \small
    \begin{tabular}{|p{3.6cm}|c|c|c|}
        \hline
        \textbf{Tiêu chí} & \textbf{\cite{jeong2026alca}} &
        \textbf{\cite{zhu2026spatiotemporal}} & \textbf{Khóa luận} \\
        \hline
        Số kho & $1$ & $4$ & $10$ \\ \hline
        Ràng buộc sức chứa & một kho & không xét & riêng từng kho \\ \hline
        Nguồn dữ liệu & tổng hợp & M5 & M5 \\ \hline
        Hàm mục tiêu & không quy về tiền & chi phí & chi phí \\ \hline
        Ghép nối giữa tác tử & không & phần thưởng chung & sức chứa dùng chung \\ \hline
        So sánh cùng mức phục vụ & không & không & có \\
        \hline
    \end{tabular}
\end{table}

\subsection{Khoảng trống nghiên cứu và định vị đóng góp}
\label{subsec:khoangTrong}

Bảng~\ref{tab:tongQuanCongTrinh} tóm tắt so sánh tổng quan giữa khóa luận này và các công trình nghiên cứu liên quan tiêu biểu.

\begin{table}[htbp]
    \centering
    \small
    \caption{So sánh tổng quan các công trình nghiên cứu liên quan}
    \label{tab:tongQuanCongTrinh}
    \setlength{\tabcolsep}{4pt}
    \renewcommand{\arraystretch}{1.25}
    \begin{tabularx}{\textwidth}{
        |>{\hsize=1.2\hsize\raggedright\arraybackslash}X
        |>{\hsize=1.2\hsize\raggedright\arraybackslash}X
        |>{\hsize=0.7\hsize\centering\arraybackslash}X
        |>{\hsize=0.9\hsize\centering\arraybackslash}X
        |>{\hsize=0.8\hsize\centering\arraybackslash}X
        |>{\hsize=1.2\hsize\centering\arraybackslash}X|
    }
        \hline
        \textbf{Nghiên cứu} & \textbf{Cấu trúc mạng} & \textbf{Số SKU} & \textbf{Chia sẻ tham số} & \textbf{Dữ liệu thực} & \textbf{Thuật toán} \\
        \hline
        Oroojlooyjadid (2022) & 4 cấp chuỗi & 1 & Không & Không & DQN \\ \hline
        Gijsbrechts (2022) & 2 cấp / đa cấp & $<10$ & Không & Không & A2C / PPO \\ \hline
        Hubbs et al. (2020) & Đa cấp lost-sales & Biến đổi & Không & Không & DQN / PPO \\ \hline
        Ding et al. (2022) & 1 cửa hàng & Nhiều SKU & Có & Không & CD-PPO / IPPO \\ \hline
        Yang et al. (2023) & Đa cấp / Đa kho & $>2000$ & Có & Có & IPPO / MAPPO \\ \hline
        Kotecha (2024) & Đa điểm phân tán & Biến đổi & Có & Không & IPPO / MAPPO \\ \hline
        Liu et al. (2025) & Đa cấp phức tạp & Nhiều SKU & Có & Có & HAPPO \\ \hline
        Khóa luận này & 10 kho song song & 30 & Có & M5 Walmart & IPPO \\
        \hline
    \end{tabularx}
\end{table}

Đóng góp của khóa luận được định vị theo hướng tái lập và mở rộng có kiểm soát đối với ứng dụng học tăng cường đa tác tử trong kiểm soát tồn kho, tập trung vào ba trọng tâm sau:
\begin{enumerate}
    \item Khảo sát cấu hình 10 kho song song thuần túy (non-multi-echelon) với ràng buộc sức chứa kho dùng chung ở quy mô $300$ cặp kho--mặt hàng, tạo điểm đối sánh trực tiếp với các mô hình đa cấp (multi-echelon) của MABIM \parencite{yang2023mabim} và Liu và cộng sự \parencite{liu2025happo}.
    \item Thực hiện khảo sát loại trừ có hệ thống nhằm định lượng tác động của tín hiệu phần thưởng phân rã cục bộ kèm ràng buộc tỷ lệ đáp ứng trên cửa sổ trượt, làm rõ cơ chế gán tín dụng và tốc độ hội tụ của IPPO ở quy mô 300 tác tử.
    \item Hiệu chỉnh môi trường mô phỏng lost-sales trực tiếp từ dữ liệu bán lẻ thực tế M5 Walmart và thực hiện kiểm chứng đối chứng có kiểm định thống kê với các chính sách cổ điển EOQ, $(s,S)$ và Newsvendor.
\end{enumerate}